\clearpage
\thispagestyle{empty}
\pdfbookmark[0]{Dioptrische Untersuchungen. Ausgabe und Quellen}{dioptrische-source-notice}
\begin{center}\Large\textbf{Dioptrische Untersuchungen}\end{center}

KI-Übertragung des deutschen Originals in modernes Deutsch mit KI-erstelltem Satz. Grundlage ist Carl Friedrich Gauss, \textit{Werke}, Band V (Göttingen, 1867), Blätter 249--282 der digitalen Vorlage. Die Ausgabe enthält das Titelblatt mit Impressum, den wiederholten Titel, die Einleitung und alle 22 Artikel. Die Arbeit wurde am 10.~Dezember 1840 vorgelegt. Das Titelblatt nennt den ersten Band der Abhandlungen der Göttinger Gesellschaft und das Göttinger Impressum von 1843.

Zum Vergleich dienten das Titelblatt, die Anmerkung des Übersetzers, die Einleitung und die Artikel 1--22 in Bernard Rosetts englischem Typoskript \textit{Optical Investigations} vom Juli 1941. Rosett verfasste auch den eigenständigen Anhang. Herangezogen wurde außerdem der Anfang des von W.~H.~Miller mitgeteilten englischen Auszugs in Richard Taylor (Hrsg.), \textit{Scientific Memoirs}, Bd.~III (1843), Art.~XIX, S.~490--498. Die Artikel \textit{Brechung} und \textit{Linsen} des Karlsruher Instituts für Technologie dienten als Quellen zur heutigen deutschen Fachsprache der Optik.

\noindent\href{https://gdz.sub.uni-goettingen.de/id/PPN236006339}{Deutsches Original (1867)}\quad
\href{https://wp.optics.arizona.edu/jsasian/wp-content/uploads/sites/33/2016/02/Gauss.pdf}{Rosett (1941)}\quad
\href{https://archive.org/details/scientificmemoir03memo/page/490/mode/1up}{Englischer Auszug (1843)}\quad
\par\smallskip\noindent\href{https://labs.physik.kit.edu/prakt-klass-physik-manuals/P1/Geometrische_Optik/doc/Hinweise-Brechung/}{KIT: Brechung}\quad
\href{https://labs.physik.kit.edu/prakt-klass-physik-manuals/P1/Geometrische_Optik/doc/Hinweise-Linsen/}{KIT: Linsen}
\clearpage
